%%%%%%%%%%%%%%%%%%%%%%%%%%%%%%%%%%%%%%%%%%%%%%%%%%%%%%%%%%%%%%%%%%%%%%%%%%%%%%%%%%%%%%%%%%%%%%%%%%%%
\chapter{Introduction} %%%%%%%%%%%%%%%%%%%%%%%%%%%%%%%%%%%%%%%%%%%%%%%%%%%%%%%%%%%%%%%%%%%%%%%%%%%%%
\label{sec:intro-motivation}
%%%%%%%%%%%%%%%%%%%%%%%%%%%%%%%%%%%%%%%%%%%%%%%%%%%%%%%%%%%%%%%%%%%%%%%%%%%%%%%%%%%%%%%%%%%%%%%%%%%%


% \todo{Write introduction paragraph.}


%%%%%%%%%%%%%%%%%%%%%%%%%%%%%%%%%%%%%%%%%%%%%%%%%%%%%%%%%%%%%%%%%%%%%%%%%%%
% \section{Motivation} %%%%%%%%%%%%%%%%%%%%%%%%%%%%%%%%%%%%%%%%%%%%%%%%%%%%
%%%%%%%%%%%%%%%%%%%%%%%%%%%%%%%%%%%%%%%%%%%%%%%%%%%%%%%%%%%%%%%%%%%%%%%%%%%


% \todo{Explain why improving error messages is important for novice programmers.}


\begin{justify}
All programmers have to start at the beginning: What is programming? What is required to write and execute code? What are variables and methods and how are they used?
From there, learners gradually encounter key concepts and algorithms, and later data structures such as maps and sets.
However, moving from basic syntax to independently writing correct programs is a long process, which also includes implementing corresponding unit tests to validate solutions.
\end{justify}

\begin{justify}
Understanding and interpreting error messages is a fundamental challenge for novice programmers.
Error messages serve as an essential form of feedback, enabling programmers to diagnose issues in their code.
In practice, standard diagnostics - especially those produced by compilers and testing frameworks - are often technical, verbose, and difficult for beginners to interpret.
\end{justify}

\begin{justify}
This becomes particularly apparent in introductory programming courses where students receive exercise templates with predefined tests and are expected to implement missing functions or data structures.
In such settings, learners are frequently unfamiliar with unit testing and are immediately confronted with stack traces and framework-specific output formats.
\end{justify}

\begin{justify}
In \fs{}, multiple testing frameworks and assertion libraries are available.
A commonly used option is \href{https://learn.microsoft.com/en-us/dotnet/api/microsoft.visualstudio.testtools.unittesting}{Microsoft.VisualStudio.TestTools.UnitTesting}.
For example, using \href{https://learn.microsoft.com/en-us/dotnet/api/microsoft.visualstudio.testtools.unittesting.assert.areequal?view=mstest-net-3.8}{\texttt{Assert.AreEqual}} can yield the following output on failure:
\end{justify}
\begin{fsharpoutput}%
{Example Ouput of \texttt{Assert.AreEqual}}%
{exampleOutputForAreEqual}
Assert.AreEqual failed. Expected:<[1; 2; 4]>. Actual:<[1; 4; 2]>.

    at MapTests.TestsMapSortedList.Tests.testSetBeispiele() in /code/TestsMapSortedList.fs:line 53
    at System.RuntimeMethodHandle.InvokeMethod(Object target, Void** arguments, Signature sig, Boolean isConstructor)
    at System.Reflection.MethodBaseInvoker.InvokeWithNoArgs(Object obj, BindingFlags invokeAttr)
\end{fsharpoutput}

\noindent For more advanced programmers, this output is usually sufficient: they can quickly identify the relevant mismatch and ignore the remaining details.
See Section \ref{sctn:errorMessages} (page \pageref{sctn:errorMessages}) for a discussion of common issues in diagnostic messages and strategies for extracting key information.
In my experience as a tutor for a first-semester lecture, I frequently observe students struggling with such messages, unsure what went wrong and how to fix their code.
Instead of guiding them toward a solution, poorly structured error messages can lead to frustration and trial-and-error debugging, which may hinder learning and delay progress on key programming concepts.\\[5pt]

Didactically, complex concepts are typically introduced gradually, often by reducing incidental complexity at the beginning:
\begin{itemize}
    \item Lightweight IDE's are often preferred to heavy, feature-rich IDE's
    \item Simplistic I/O is introduced before full input validation only valid inputs and error handling (e.g., malformed inputs, exceptions, and missing values).
\end{itemize}

\noindent The same  Diagnostic verbosity can be staged for testing: Initially showing only the most relevant information and revealing additional technical details as students gain experience.
This thesis explores how a similar scaffolding approach can be applied to testing feedback.
By providing more minimalistic and structured output, the goal is to help students interpret failures more effectively and support a more efficient and engaging learning process.


%%%%%%%%%%%%%%%%%%%%%%%%%%%%%%%%%%%%%%%%%%%%%%%%%%%%%%%%%%%%%%%%%%%%%%%%%%%
\section{Research Problem and Objectives} %%%%%%%%%%%%%%%%%%%%%%%%%%%%%%%%%
\label{sec:intro-research-problem}
%%%%%%%%%%%%%%%%%%%%%%%%%%%%%%%%%%%%%%%%%%%%%%%%%%%%%%%%%%%%%%%%%%%%%%%%%%%


% \todo{Define the problem of unclear error messages and the goals of the thesis.}

The primary problem addressed in this thesis is that, in introductory programming settings,
the feedback produced by common testing frameworks is often not presented in a form that is immediately interpretable for novice programmers.
In particular, test failures frequently mix the core mismatch information (what was expected versus what occurred)
with framework-specific detail such as stack traces and internal method calls.
As a consequence, beginners may struggle to identify the relevant part of the output and to translate it into concrete corrective actions in their own code.

\noindent From a technical perspective, this raises the question of how test feedback can be reformatted in a way that is
both \textbf{(i)} more consistent and cognitively simpler for first-semester students and \textbf{(ii)} still compatible with established testing infrastructures.
This thesis therefore focuses on the design and implementation of a lightweight assertion and reporting layer, \emph{Assertify}, that produces a compact, structured summary of test failures.
The objectives of this thesis are as follows:
\begin{itemize}[nosep]
    \item Analyse the structure and typical content of failing test output in beginner exercises, with a focus on which parts are essential for interpreting the failure.
    \item Design and implement a custom test/assertion library that enforces a consistent reporting scheme (e.g., a compact \texttt{expected = actual} summary) on top of existing frameworks.
    \item Reflect on design and implementation decisions by identifying current limitations of the library and outlining technically feasible extensions and future improvements.
\end{itemize}

To achieve these objectives, this thesis presents the implementation of \textit{Assertify} as a wrapper around existing mechanisms for unit and property-based testing,
and discusses how the reporting format can be further generalised, extended, and integrated into teaching workflows.


%%%%%%%%%%%%%%%%%%%%%%%%%%%%%%%%%%%%%%%%%%%%%%%%%%%%%%%%%%%%%%%%%%%%%%%%%%%
\section{Research Questions} %%%%%%%%%%%%%%%%%%%%%%%%%%%%%%%%%%%%%%%%%%%%%%
%%%%%%%%%%%%%%%%%%%%%%%%%%%%%%%%%%%%%%%%%%%%%%%%%%%%%%%%%%%%%%%%%%%%%%%%%%%


% \todo{List the key questions this thesis seeks to answer.}


This thesis seeks to answer the following key research questions:
\begin{itemize}
    \item \textbf{RQ1 (Problem characterisation):}
        \begin{justify}Which structural characteristics of compiler and testing-framework outputs make failing tests difficult to interpret for novice programmers in introductory programming exercises?\end{justify}
    \item \textbf{RQ2 (Design requirements):}
        \begin{justify}Which requirements can be derived for novice-oriented test feedback with respect to consistency, cognitive load, and the ``actionability'' of failure information?\end{justify}
    \item \textbf{RQ3 (Design proposal):}
        \begin{justify}How can a uniform ``tested expression / expected / actual'' reporting scheme be defined such that it remains applicable across different kinds of tests (example-based unit tests and property-based tests)?\end{justify}
    \item \textbf{RQ4 (Technical feasibility):}
        \begin{justify}How can such a reporting scheme be implemented in \fs{} using Quotations and Unquote for expression capture and evaluation, while remaining compatible with existing tooling through wrappers around\\\texttt{TestTools.UnitTesting} and \texttt{FsCheck}?\end{justify}
    \item \textbf{RQ5 (Property-based integration):}
        \begin{justify}How can property-based testing outcomes---including generated counterexamples and shrinking results---be represented within the same output format using \texttt{Checkify.CheckTest} and generators?\end{justify}
    \item \textbf{RQ6 (Reflection and future work):}
        \begin{justify}What are the current limitations of the implemented library (e.g., supported assertion forms, expression coverage, formatting constraints, integration boundaries), and which extensions are technically feasible for future versions?\end{justify}
\end{itemize}

\noindent By addressing these questions, this thesis positions its contribution as the design and implementation of a test feedback layer that simplifies and
standardises failure reporting for novice programmers, and as a critical reflection on how such a library can be extended and integrated into teaching contexts.


%%%%%%%%%%%%%%%%%%%%%%%%%%%%%%%%%%%%%%%%%%%%%%%%%%%%%%%%%%%%%%%%%%%%%%%%%%%
\section{Thesis Structure} %%%%%%%%%%%%%%%%%%%%%%%%%%%%%%%%%%%%%%%%%%%%%%%%
\label{sec:intro-structure}
%%%%%%%%%%%%%%%%%%%%%%%%%%%%%%%%%%%%%%%%%%%%%%%%%%%%%%%%%%%%%%%%%%%%%%%%%%%


% \todo{Briefly outline the chapters and their purpose.}



This thesis is structured as follows:
\begin{itemize}
    \item \textbf{Chapter 1:} Introduction
    \item \textbf{Chapter 2:} Related Work
    \item \textbf{Chapter 3:} Assertify/Checkify
    \item \textbf{Chapter 4:} Evaluation and Results
    \item \textbf{Chapter 5:} Discussion

    \item \textbf{Chapter 2: Related Work}
        \begin{justify}Reviews research on novice programmers' difficulties with compiler and test diagnostics, and summarises prior approaches to improving the structure and usefulness of error feedback.
            The chapter establishes the conceptual motivation for presenting failures in a more consistent and cognitively simple format.\end{justify}

    \item \textbf{Chapter 3: Analysis of Previous Semester Data}
        \begin{justify}Analyses artefacts from earlier course iterations (e.g., submissions, test outcomes, and common failure patterns) to identify recurring sources of errors and
            to contextualise which kinds of feedback are most frequently encountered by students.
            This chapter motivates the design goals of a novice-oriented reporting scheme by grounding the problem in a concrete teaching setting.\end{justify}

    \item \textbf{Chapter 4: Problem Analysis}
        \begin{justify}Examines typical compiler and testing-framework outputs as they appear in beginner exercises and characterises their technical properties
            (e.g., verbosity, inconsistency, mixing of essential information with framework internals). Based on this analysis,
            the chapter derives design requirements for improved test failure reporting, focusing on a stable ``tested expression / expected / actual'' summary as the primary interface.\end{justify}

    \item \textbf{Chapter 5: Assertify --- Design and Implementation}
        \begin{justify}Presents the design rationale and implementation of \emph{Assertify} as a lightweight test and assertion library that wraps existing testing infrastructures.
            It explains how a consistent \texttt{expected = actual} output format is produced, including the use of Quotations and Unquote for expression capture and evaluation,
            as well as wrappers around \texttt{UnitTesting} and \texttt{FsCheck}. The chapter also discusses implementation choices and their implications for extensibility.\end{justify}

    \item \textbf{Chapter 6: Evaluation and Results}
        \begin{justify}Evaluates Assertify primarily from a technical and qualitative perspective, focusing on the properties the implementation achieves
            (e.g., consistency of output, coverage of supported assertion forms, integration with unit and property-based testing workflows).
            Where applicable, the chapter contrasts Assertify output with baseline framework output to illustrate differences in interpretability,
            while clearly separating observations from generalisable claims.\end{justify}

    \item \textbf{Chapter 7: Discussion and Future Work}
        \begin{justify}Reflects on the contribution and limitations of Assertify and outlines potential extensions.
            This includes incremental improvements within the wrapper approach (e.g., broader expression coverage, richer formatting layers, IDE integration),
            as well as a discussion of what would become possible with more extensive redesign efforts, such as implementing a dedicated testing framework from scratch instead of relying on wrappers.
            The chapter concludes by positioning Assertify as a step toward more pedagogically aligned testing feedback and summarising open directions for future work.\end{justify}

    \item \textbf{Chapter 8: Conclusion}
        \begin{justify}Summarises the thesis contributions, revisits the research questions, and highlights the main insights regarding the feasibility and design space of novice-oriented test feedback.\end{justify}
\end{itemize}

\noindent nnnnnnnThis structure emphasizes the design and implementation contribution of the thesis and concludes with a focused discussion of technical possibilities and future extensions,
ranging from incremental refinements of the current library to more fundamental redesigns of testing feedback infrastructures.


\begin{comment}
This thesis is structured as follows:
\begin{itemize}
    \item \textbf{Chapter 2: Related Work} -- Reviews existing research on challenges faced by novice programmers, the role of error messages in learning, and previous attempts at improving error feedback.
    \item \textbf{Chapter 3: Analysis of Previous Semester Data} -- Presents an analysis of past student submissions, including compilation success rates and test pass rates, to identify common issues.
    \item \textbf{Chapter 4: Problem Analysis} -- Examines specific problems with existing error messages, including real-world examples and gaps in the information they provide.
    \item \textbf{Chapter 5: Assertify --- A Library for Improving Error Messages} -- Introduces and explains the implementation of Assertify, a tool designed to enhance error feedback.
    \item \textbf{Chapter 6: Further Methodology (VS Code Extension?)} -- TODO
    \item \textbf{Chapter 7: Evaluation and Results} -- Presents an analysis of the impact of modified error messages, including both quantitative and qualitative findings.
    \item \textbf{Chapter 8: Conclusion} -- Summarizes key findings, discusses their implications, and provides recommendations for future work.
\end{itemize}
This structure ensures a logical flow, moving from problem identification to solution development and evaluation, ultimately contributing to a better understanding of how error messages affect novice programmers.




ACTUAL = EXPECTED           ->              f(x) = 0        oder        test 1 2 = 3 (mehr natürlich zu lesen)


- Intro
    Tests = bad -> kann man die vielleicht besser machen?
- Related Work
    Da wollt ich nochmal nachfragen was für Literatur Themen du jetzt genau wichtig findest im aktuellen Kontext der Arbeit
- Problem Analysis
    Hier würde ich eben die konkreten Schwierigkeiten der Fehlermeldungen erklären und wie daraus dieses `actual = expected` enstand
    und warum man Quotations braucht?
- Assertify
    - Die konkreten Details zu Assertify/Checkify
    - CaptureRunner
    - eval, apply, etc.
- Evaluation and Results
    Joa :D
- Discussion
    Und hier dann eben wie man es hätte besser machen können und wie man das ganze erweitern oder besser, besser machen könnte
    Mit `FSharp.Compiler` kann man ja auch programmatisch Code ausführen und Code auswerten und dann mit der diagnostics arbeiten.
    Das könnte man halt nutzen um spezieller auf Fehlersuche zu gehen im Studicode.
    Man bekommt dann auch die genaue stelle im Code. Mit FParsec könnte man an der Stelle anfangen custom parser drauf zu klatschen mit immer größerem Radius,
    die nach Dingen suchen wie `9` statt `9N` und andere der häufigsten Fehler.
    Wenn man dann nämlich den ErrorCode und die Message der diagnostics hinzunimmt kann bestimmt viel anfangen.

Ziel: Lesbarere Error messages: "Expected : Got :"
Problem:
    Bestehende Property-based-tests der Form
        ∀ a ∈ A -> ∀ y ∈ B(a) -> f (x,y) = S.f(x,y),
    wollen dass bestehende FSCheck tests relativ einfach geported werden können.
Solution Teil 1:
    Quote / Unquote:
        Erlaubt String-Darstellung zur Laufzeit von Ausdrücken,
        insb. der Form "f x = S.f y", für x,y die zur Laufzeit bestimmt werden.
        → Insb für den Fall fun a -> fun b -> foo a b = bar a b geht das in Check mit quotation um diesen Rumpf und
        dann Expressins.apply und Expressions.eval.
        → für den Fall ∀ a ∈ A -> ∀ y ∈ B(a) -> f (x,y) = S.f(x,y) mit CheckTest über die Trennung des
        "lambda prefixes" ∀ a ∈ A -> ∀ y ∈ B(a) -> und der "expression under test" f (x,y) = S.f(x,y),
        was einen kleinen Overhead im Vergleich zu einem reinen FSCheck Test darstellt.
Solution Teil 2:
    Eingreifen in die Ausführung von FSCheck über CaptureRunner um die Darstellung der Fehlermeldung in unserem Format zu erreichen.


\end{comment}


%%%%%%%%%%%%%%%%%%%%%%%%%%%%%%%%%%%%%%%%%%%%%%%%%%%%%%%%%%%%%%%%%%%%%%%%%%%%%%%%%%%%%%%%%%%%%%%%%%%%
% EOF %%%%%%%%%%%%%%%%%%%%%%%%%%%%%%%%%%%%%%%%%%%%%%%%%%%%%%%%%%%%%%%%%%%%%%%%%%%%%%%%%%%%%%%%%%%%%%
%%%%%%%%%%%%%%%%%%%%%%%%%%%%%%%%%%%%%%%%%%%%%%%%%%%%%%%%%%%%%%%%%%%%%%%%%%%%%%%%%%%%%%%%%%%%%%%%%%%%